% Reviewer editorial proposal for r09; not adopted or inserted into a manuscript.
% Replace the transition after Theorem thm:smallbatch and before cor:eventual.
% Keep the existing bracketing explanation below; do not also retain the old
% sentence "The privacy variance alone eventually exceeds ... giving:".
% This replacement targets the original transition in ORIGINAL_TRANSITION.tex.
% If merging with the separate offset proposal, retain its calibration-identity
% and horizon-substitution explanation; this clarification does not replace it.
% Existing labels and fixed-parameter hypotheses are reused, not redefined.
The exact profile limit is inverted by bracketing the calibrated noise between
$\Delta/(\sqrt{2\log T}+s_{E,\delta}\pm\eta)$ for each $\eta>0$.
The general profile limit and this inversion are proved in the appendix.

The divergence in~\eqref{eq:privacydivergence} concerns the Gaussian
variance of the $E$-normalized estimator, not the noise variance of one
release. Since the raw noise is recalibrated at every horizon, the factor
$T$ alone does not imply that $v_{E,T}=T\sigma_{E,T}^2/E^2$ diverges.
Under the fixed-parameter hypotheses above, the theorem gives
$\sigma_{E,T}\sim\Delta/\sqrt{2\log T}$ and hence
$v_{E,T}\sim\Delta^2T/(2E^2\log T)\to\infty$.
Thus, for sufficiently large transcript horizons, the privacy variance alone
exceeds the fixed full-batch variance $v_{\mathrm{FB}}$; the nonnegative
sampling term makes the same comparison hold for every $S\geq0$.
This is an eventual comparison, not an assertion of optimality at every finite
horizon. If only the final average is disclosed and calibration is performed
for that output instead, Proposition~\ref{prop:average-uniform-bound} provides
a uniform upper bound. That bound is neither a limiting formula nor a privacy
guarantee for additionally disclosed individual releases.

The resulting comparison is stated in the following corollary.
